\begin{titlepage}
\begin{center}
NATIONAL ECONOMICS UNIVERSITY, HANOI \\
COLLEGE OF TECHNOLOGY \\
FACULTY OF DATA SCIENCE AND ARTIFICIAL INTELLIGENCE
\end{center}

\vspace{0.4cm}

\begin{center}
\includegraphics[width=3.4cm]{Images/NEU logo.png}
\end{center}

\vspace{0.4cm}

\begin{center}
\begin{tabular}{c}
\multicolumn{1}{c}{\textbf{{\LARGE CREDIT CARD FRAUD DETECTION}}}\\
\\[0.3cm]
\hline
\\[0.3cm]
\multicolumn{1}{c}{\textbf{{\Large FINAL REPORT}}}\\
\\[0.2cm]
\textbf{{\Huge Topic: Credit Card Fraud Detection}}\\[0.2cm]
\textbf{{\large Handling Class Imbalance and Designing a Cost-Aware Operating Policy}}\\
\\[0.3cm]
\hline
\end{tabular}
\end{center}

\vspace{0.9cm}

\begin{center}
\begin{tabular}{rrll}

\hspace{5 cm} & Instructor: Ph.D. Nguyễn Trọng Nghĩa  & &\\

& Student: &

    Nguyễn Tuấn Sỹ & 11230775 \\
& & Nguyễn Tiến Công & 11230739 \\
& & Nguyễn Nho Thắng & 11230778 \\
& & Nguyễn Quang Cường & 11230740 \\
& & Nguyễn Gia Phong & 11230768 \\

\end{tabular}
\end{center}

\vspace{0.45cm}

\begin{center}
\begin{tabular}{p{12.4cm}}
\textbf{\textit{Abstract.}} Fraud detection on transactional data is an
imbalanced, cost-sensitive problem: only 0.167\% of the 283{,}726 transactions
in the ULB credit card dataset are fraudulent, so accuracy is uninformative
and the operating point matters more than the classifier. We evaluate 27
configurations spanning five model families (Logistic Regression, Decision
Tree, Random Forest, XGBoost, Isolation Forest) and five imbalance-handling
techniques under a leakage-controlled workflow, and we translate the winning
model into a deployable operating policy. Random Forest with SMOTE at a 1:100
ratio is selected (F1 = 0.8605, Average Precision = 0.8239), reaching 78.17\%
recall at a 95\% precision floor, or 73.24\% with a leakage-free threshold.
Central to the business answer is the capacity constraint: at 141{,}863
transactions per day no single threshold can satisfy both the 95\% precision
floor and a 100-review-per-day analyst capacity --- it either demands 1.9
analysts or throttles recall to 40.9\%. A two-tier policy that auto-blocks
above 0.810 and reviews only the 0.175--0.810 band resolves this, attaining
81.0\% recall at USD 154.65 per 85{,}118 transactions with just 70 reviews per
day --- an operational gain over the cheapest single threshold near that
operating point, which needs 232 reviews per day for comparable recall. We also quantify the effect of tuning the
threshold on the test set, report an observation consistent with a wider hyper-parameter
search overfitting the cross-validation folds, and demonstrate that permutation
importance and SHAP both rank \texttt{Amount} first where the default impurity
importance ranks it 21st of 35.

\vspace{0.25cm}
\textbf{\textit{Keywords:}} credit card fraud detection; class imbalance;
SMOTE; cost-sensitive learning; precision--recall trade-off; threshold
calibration; operating policy; SHAP.
\end{tabular}
\end{center}

\vspace{0.45cm}
\begin{center}
\begin{tabular}{p{11.5cm}}
\textit{``All quantitative results reported in this document were produced by
the authors' own pipeline; the random seed is fixed at 42, so every figure
in this report is exactly reproducible. The source code, the six executed
Jupyter notebooks, all 44 result tables and the LaTeX source of this report
are publicly available at
\href{https://github.com/eqflashheqr-cmyk/DataMining}{\texttt{https://github.com/eqflashheqr-cmyk/DataMining}}
under the MIT licence. The raw dataset (public on Kaggle) and the trained
\texttt{.joblib} models are excluded for size; every artefact is regenerated
by the scripts in \texttt{src/}.''}
\end{tabular}
\end{center}

\vfill

\begin{center}
{\footnotesize Hanoi, October 2026}
\end{center}

\end{titlepage}

% The titlepage environment resets the page counter, which made the footer read
% "Page 19/19" on a 20-page PDF. Restore the physical numbering so that the
% printed page number always matches the page count.
\setcounter{page}{2}

\thispagestyle{empty}

\tableofcontents
